\subsection{Exciton Formation}

When an electron is promoted from the valence band to the conduction band, for example by optical absorption, it leaves behind a hole of effective charge $+e$. Electrons in the conduction band can feel the potential created by the holes in the valence band. This electrostatic attraction can be modeled by a Coulomb potential, which allows us to solve the Schrödinger equation for a composite particle called an \textbf{exciton}, made of a conduction band electron and a valence band hole.

This quasiparticle has a reduced mass defined by $\frac{1}{\mu} = \frac{1}{m_e} + \frac{1}{m_h}$ and a quantized energy spectrum of the hydrogen-like type.

In the absence of Coulomb interaction, the minimum energy required to create a free electron-hole pair corresponds to the gap energy $E_g$. However, because of the attractive Coulomb potential between the electron and the hole, they can form a bound state. The accessible energy levels then sit below the conduction band. These states can be populated directly through the process described above, but also through nonradiative relaxation mechanisms such as the Auger effect or exciton-phonon interactions. As pointed out by Mao et al.~\cite{mao2004}, this formation process is very fast, often under $1 \text{ ps}$ for wide bandgap materials.

The relative motion of this electron-hole pair can be modeled by a hydrogen-like Schrödinger equation. The stationary equation for the relative coordinate $\mathbf{r}$ reads:
\begin{equation}
    \left[ \frac{\hbar^2}{2\mu} \nabla^2 + \frac{e^2}{4\pi \epsilon_0 \epsilon_r r} \right] \psi(\mathbf{r}) = E_b \psi(\mathbf{r}),
\end{equation}
where $\epsilon = \epsilon_0 \epsilon_r$ is the dielectric permittivity of the medium and $E_b > 0$ is the exciton binding energy.

The relation between the gap energy $E_g$, the binding energy term, and the excitation energy $E_n$ required for the $n$-th exciton state is given by:
\begin{equation}
    E_g = E_n + \frac{R_\chi}{n^2},
\end{equation}
where $R_\chi$ is the effective exciton Rydberg energy:
\begin{equation}
    R_\chi = \frac{\mu e^4}{2(4\pi \epsilon)^2 \hbar^2} = R_H \left(\frac{\mu}{m_0}\right) \frac{1}{\epsilon_r^2}.
\end{equation}
Here, $R_H \approx 13.6\text{ eV}$ is the Rydberg constant of the hydrogen atom, $m_0$ is the free electron mass, and $\epsilon_r$ is the relative dielectric constant of the material.

To evaluate the spatial extent of the exciton wavefunction, we define the exciton Bohr radius $a_\chi$:
\begin{equation}
    a_\chi = \frac{4\pi \epsilon \hbar^2}{\mu e^2} = a_0 \cdot \epsilon_r \left(\frac{m_0}{\mu}\right),
\end{equation}
where $a_0 \approx 0.53\text{ \AA}$ is the atomic Bohr radius.

In wide bandgap materials, electrons are strongly bound to their atoms, which means the relative dielectric permittivity is low for $\mathrm{SiO}_2$ or alkali halides, where the relative dielectric constant is typically of order $2$ to $4$. As a result, the material provides weak dielectric screening of the Coulomb attraction, maintaining a deep potential well.

A low permittivity directly leads to a reduction of the exciton Bohr radius, down to a few angstroms. When $a_\chi$ becomes comparable to the interatomic spacing of the lattice, the exciton localization no longer extends over several lattice atoms but shifts to a strongly bound regime. The binding energy also increases, preventing thermal dissociation and favoring ultrafast energy relaxation. During this process, the excess electrostatic potential energy of the unscreened exciton is abruptly discharged into local lattice vibrations (phonons) through massive exciton-phonon coupling within a few picoseconds. This induces a local lattice distortion and leads to the formation of a self-trapped exciton (STE).
